\section{Mode A5 -- Préparation pour remise en route après défaillance}
\label{secModeA5}

\subsection{Conditions d'entrée}

Le mode A5 est atteint depuis le mode D2, une fois que l'installation a
retrouvé une configuration saine.

\subsection{Comportement attendu}

L'objectif de ce mode est de préparer une remise en route sûre de
l'installation après une défaillance, en tenant compte du fait qu'une
pièce pouvait être en cours de manipulation au moment où la défaillance
est survenue.

\begin{itemize}
    \item \textbf{Robot} : le robot devra déposer l'éventuel gobelet, s'il
    en porte un.
    \item \textbf{Convoyeur / guichet} : l'accès au guichet reste interdit
    pendant toute la durée de ce mode.
    \item \textbf{Opérateur} : dans la mesure où l'installation ne peut
    pas, à elle seule, s'assurer qu'aucune pièce n'est restée engagée dans
    l'une des aires de traitement (cuves), c'est à l'opérateur qu'il
    revient :
    \begin{itemize}
        \item d'autoriser la dépose de l'éventuelle pièce restée saisie par
        le robot ;
        \item de confirmer que l'ensemble des aires de traitement sont
        effectivement vides.
    \end{itemize}
    \item \textbf{IHM} : la situation doit être signalée à l'opérateur, et
    les deux confirmations attendues de sa part (autorisation de dépose,
    aires de traitement vides) doivent lui être demandées explicitement.
\end{itemize}

\begin{UPSTIwarning}{Sécurité de la remise en route}
Ce mode ne doit pas être quitté tant que l'opérateur n'a pas explicitement
confirmé les deux conditions attendues de lui. Cette validation humaine est
indispensable dès lors que l'installation ne dispose pas des moyens de
détection nécessaires pour vérifier elle-même l'état des aires de
traitement et du système de préhension.
\end{UPSTIwarning}

\subsection{Conditions de sortie}

\begin{itemize}
    \item Vers le mode \textbf{A6} (\emph{Mise en conditions initiales}),
    une fois que l'opérateur a autorisé la dépose de l'éventuelle pièce
    restée saisie et confirmé que toutes les aires de traitement sont
    vides.
\end{itemize}
